\section{Notation and formal conventions}
\label{app:notation}

\begin{longtable}{@{}p{0.22\textwidth}p{0.70\textwidth}@{}}
\toprule
Symbol & Meaning \\
\midrule
\endfirsthead
\toprule
Symbol & Meaning \\
\midrule
\endhead
\(\iPRA\) & Intuitionistic primitive recursive arithmetic. \\
\(\PRAAx\) & The axiom family used by the embedded ordinary-iPRA calculus in
\code{ConciseIPRA}. \\
\(\Gamma\vdash_{\PRAAx}A\) & An embedded derivation, represented in Lean by
\code{Prf PRAAx Gamma A}. \\
\(\PTm n\), \(\Fm\,\PTm\,n\) & Object terms and formulas with \(n\) free variables. \\
\(\chi_A\) & Primitive-recursive characteristic term of a bounded formula \(A\). \\
\(r_F(d)\) & Largest input row read while computing output demand \(d\). \\
\(j_F(d)\) & Largest input contract row used in the proof of output demand \(d\). \\
\(W_i(a)\) & Backward viability bit for candidate \(a\) at level \(i\). \\
\(\delta_k\) & Same-level diameter bound among viable observations. \\
\(\eta_k\) & Adjacent-level motion bound along a compatibility edge. \\
\(\Gamma_G\) & Permanent boundary context of an \(\FRNF\) graph. \\
\bottomrule
\end{longtable}

Natural numbers are the object domain.  Integers, rationals, finite sequences, vectors,
matrices, and finite graphs are coded by naturals.  Rational inequalities are reduced
to signed-integer or natural-number equations after clearing positive denominators.
``Primitive recursive'' always means first-order primitive recursive over \(\N\), not
an arbitrary higher-type functional.

All generic constructions are schemas over fixed syntax.  When a program or transition
is described as a parameter in the metatheory, the object theorem uses the concrete
primitive-recursive symbol produced by static specialization.
